\documentclass[10pt,a4paper]{amsart}
\usepackage{amsmath,amssymb,amsthm,graphicx}
% Retained preamble for the Stacks Project chapter snapshot in this directory.
% The chapter's own preamble.tex is a copy of the shared Stacks Project preamble; the block
% below is sliced MECHANICALLY, line for line, from it. Nothing was added, removed or
% redefined; the wrapper class, the xr-hyper/hyperref wiring and the externaldocument list
% are omitted because the assembler supplies the document class.
%
% lines taken from preamble.tex: 3, 5, 9, 11, 12, 156, 157, 158, ... (51 lines total)
%
% package coverage restored from the chapter's own preamble (2026-10-13): xy, xy-2cell

\usepackage[all]{xy}\usepackage{verbatim}
\xyoption{2cell}
\UseAllTwocells

\newenvironment{reference}{\comment}{\endcomment}
\newenvironment{slogan}{\comment}{\endcomment}
\newenvironment{history}{\comment}{\endcomment}
\theoremstyle{plain}
\newtheorem{theorem}[subsection]{Theorem}
\newtheorem{proposition}[subsection]{Proposition}
\newtheorem{lemma}[subsection]{Lemma}
\theoremstyle{definition}
\newtheorem{definition}[subsection]{Definition}
\newtheorem{example}[subsection]{Example}
\newtheorem{exercise}[subsection]{Exercise}
\newtheorem{situation}[subsection]{Situation}
\theoremstyle{remark}
\newtheorem{remark}[subsection]{Remark}
\newtheorem{remarks}[subsection]{Remarks}
\numberwithin{equation}{subsection}
\def\lim{\mathop{\mathrm{lim}}\nolimits}
\def\colim{\mathop{\mathrm{colim}}\nolimits}
\def\Spec{\mathop{\mathrm{Spec}}}
\def\Hom{\mathop{\mathrm{Hom}}\nolimits}
\def\Ext{\mathop{\mathrm{Ext}}\nolimits}
\def\SheafHom{\mathop{\mathcal{H}\!\mathit{om}}\nolimits}
\def\SheafExt{\mathop{\mathcal{E}\!\mathit{xt}}\nolimits}
\def\Sch{\mathit{Sch}}
\def\Mor{\mathop{\mathrm{Mor}}\nolimits}
\def\Ob{\mathop{\mathrm{Ob}}\nolimits}
\def\Sh{\mathop{\mathit{Sh}}\nolimits}
\def\NL{\mathop{N\!L}\nolimits}
\def\CH{\mathop{\mathrm{CH}}\nolimits}
\def\proetale{{pro\text{-}\acute{e}tale}}
\def\etale{{\acute{e}tale}}
\def\QCoh{\mathit{QCoh}}
\def\Ker{\mathop{\mathrm{Ker}}}
\def\Im{\mathop{\mathrm{Im}}}
\def\Coker{\mathop{\mathrm{Coker}}}
\def\Coim{\mathop{\mathrm{Coim}}}
\DeclareMathSymbol{\boxtimes}{\mathbin}{AMSa}{"02}
\def\QCohstack{\mathcal{QC}\!\mathit{oh}}
\def\Cohstack{\mathcal{C}\!\mathit{oh}}
\def\Spacesstack{\mathcal{S}\!\mathit{paces}}
\def\Quotfunctor{\mathrm{Quot}}
\def\Hilbfunctor{\mathrm{Hilb}}
\def\Curvesstack{\mathcal{C}\!\mathit{urves}}
\def\Polarizedstack{\mathcal{P}\!\mathit{olarized}}
\def\Complexesstack{\mathcal{C}\!\mathit{omplexes}}
\def\Pic{\mathop{\mathrm{Pic}}\nolimits}
\def\Picardstack{\mathcal{P}\!\mathit{ic}}
\def\Picardfunctor{\mathrm{Pic}}
\def\Deformationcategory{\mathcal{D}\!\mathit{ef}}
\usepackage{cleveref}
\crefname{originallemma}{original lemma}{original lemmas}
\Crefname{originallemma}{Original lemma}{Original lemmas}
\makeatletter
\newlabel{spaces-cohomology-lemma-check-proper-dvr}{{0ARK}{source}{}{sourcepointer-1}{}}
\newlabel{spaces-cohomology-lemma-check-proper-dvr@cref}{{[originallemma][0][]0ARK}{[1][0][]source}{}{}{}}
\makeatother
\crefname{originallemma}{original lemma}{original lemmas}
\Crefname{originallemma}{Original lemma}{Original lemmas}
\makeatletter
\newlabel{stacks-morphisms-definition-fill-in-diagram}{{0CLA}{source}{}{sourcepointer-2}{}}
\newlabel{stacks-morphisms-definition-fill-in-diagram@cref}{{[originallemma][0][]0CLA}{[1][0][]source}{}{}{}}
\makeatother
\crefname{originallemma}{original lemma}{original lemmas}
\Crefname{originallemma}{Original lemma}{Original lemmas}
\makeatletter
\newlabel{stacks-morphisms-lemma-cat-dotted-arrows-base-change}{{0CLE}{source}{}{sourcepointer-3}{}}
\newlabel{stacks-morphisms-lemma-cat-dotted-arrows-base-change@cref}{{[originallemma][0][]0CLE}{[1][0][]source}{}{}{}}
\makeatother
\crefname{originallemma}{original lemma}{original lemmas}
\Crefname{originallemma}{Original lemma}{Original lemmas}
\makeatletter
\newlabel{stacks-morphisms-lemma-helper-diagonal}{{0E8L}{source}{}{sourcepointer-4}{}}
\newlabel{stacks-morphisms-lemma-helper-diagonal@cref}{{[originallemma][0][]0E8L}{[1][0][]source}{}{}{}}
\makeatother
\crefname{originallemma}{original lemma}{original lemmas}
\Crefname{originallemma}{Original lemma}{Original lemmas}
\makeatletter
\newlabel{stacks-morphisms-lemma-locally-finite-type-locally-noetherian}{{06R6}{source}{}{sourcepointer-5}{}}
\newlabel{stacks-morphisms-lemma-locally-finite-type-locally-noetherian@cref}{{[originallemma][0][]06R6}{[1][0][]source}{}{}{}}
\makeatother
\crefname{originallemma}{original lemma}{original lemmas}
\Crefname{originallemma}{Original lemma}{Original lemmas}
\makeatletter
\newlabel{stacks-morphisms-lemma-properties-diagonal}{{04XS}{source}{}{sourcepointer-6}{}}
\newlabel{stacks-morphisms-lemma-properties-diagonal@cref}{{[originallemma][0][]04XS}{[1][0][]source}{}{}{}}
\makeatother
\crefname{originallemma}{original lemma}{original lemmas}
\Crefname{originallemma}{Original lemma}{Original lemmas}
\makeatletter
\newlabel{stacks-properties-lemma-check-property-covering}{{04XD}{source}{}{sourcepointer-7}{}}
\newlabel{stacks-properties-lemma-check-property-covering@cref}{{[originallemma][0][]04XD}{[1][0][]source}{}{}{}}
\makeatother
\title{DVR fill-in criterion for separatedness}
\author{The Stacks Project Authors}
\date{stacks-project:stacks-more-morphisms.tex:local-snapshot-20261012; GNU FDL 1.2 or later}
\begin{document}
\maketitle
\noindent\textit{Editorial note.} Statement span is lines 2840-2860 of sources/stacks-stacks-more-morphisms/pinned-original-source.tex; proof span is lines 2862-2906, the complete author proof (43 body lines).\par\medskip
\noindent\textit{Editorial note.} Navigation stubs: 7 cited-but-not-included labels are supplied as non-load-bearing pointers with their pinned tags where the snapshot tag index has them.\par\medskip
\noindent\textit{Editorial note.} extra\_packages is empty on purpose: the retained preamble already loads xy with options and a second declaration clashes in the assembler wrapper.\par\medskip
\noindent\textit{Editorial note.} Text is retained verbatim, including the author's own typos and placeholders.\par\medskip
\noindent\textit{Editorial note.} Queue index 29 of queue\_stacks-stacks-more-morphisms.json (pb=1909, stmt 2840-2860, proof 2862-2906).\par\medskip
\noindent\textit{Editorial note.} stacks-stacks-more-morphisms wave smm01 (2026-10-14): fresh adjudication accept per lane-s3/smm\_judgments.md; pool per Chief 10:5x open-pool ruling (net 59 with-proof rows, bloodline lane-s3 by anchor audit, v2 true-key namespace stacks:stacks-more-morphisms.tex, row lemma-deform-property-fp-over-ft avoided on cross-chapter key collision); 60 with-proof rows judged 45 accept / 14 reject / 1 avoid; patch-sd-envs contract spans scanned.\par\medskip
\noindent\textit{Editorial note.} Cross-chapter same-label warning (NOT an occupancy conflict, Chief ruling 2026-10-13 evening): label 'lemma-check-separated-dvr' also exists in stacks-spaces-cohomology — a different proposition in a different chapter file.\par\medskip
\noindent\textit{Allowed standard prerequisites (editorial).}\begin{itemize}
\item Explicit prerequisite boundary. The cited inputs are prerequisites; the argument in the retained span is not replaced by any cited result.
\end{itemize}\medskip
\ifdefined\hypertarget\hypertarget{sourcepointer-1}{}\fi
\noindent\textit{Source reference.} Original lemma 0ARK: sources/stacks-stacks-more-morphisms/pinned-original-source.tex, source lines 2906--2906. Cited by the retained statement or proof but not part of the collected spans; supplied only as a navigation pointer. The argument does not depend on its content.\par\medskip
\ifdefined\hypertarget\hypertarget{sourcepointer-2}{}\fi
\noindent\textit{Source reference.} Original lemma 0CLA: sources/stacks-stacks-more-morphisms/pinned-original-source.tex, source lines 2906--2906. Cited by the retained statement or proof but not part of the collected spans; supplied only as a navigation pointer. The argument does not depend on its content.\par\medskip
\ifdefined\hypertarget\hypertarget{sourcepointer-3}{}\fi
\noindent\textit{Source reference.} Original lemma 0CLE: sources/stacks-stacks-more-morphisms/pinned-original-source.tex, source lines 2906--2906. Cited by the retained statement or proof but not part of the collected spans; supplied only as a navigation pointer. The argument does not depend on its content.\par\medskip
\ifdefined\hypertarget\hypertarget{sourcepointer-4}{}\fi
\noindent\textit{Source reference.} Original lemma 0E8L: sources/stacks-stacks-more-morphisms/pinned-original-source.tex, source lines 2906--2906. Cited by the retained statement or proof but not part of the collected spans; supplied only as a navigation pointer. The argument does not depend on its content.\par\medskip
\ifdefined\hypertarget\hypertarget{sourcepointer-5}{}\fi
\noindent\textit{Source reference.} Original lemma 06R6: sources/stacks-stacks-more-morphisms/pinned-original-source.tex, source lines 2906--2906. Cited by the retained statement or proof but not part of the collected spans; supplied only as a navigation pointer. The argument does not depend on its content.\par\medskip
\ifdefined\hypertarget\hypertarget{sourcepointer-6}{}\fi
\noindent\textit{Source reference.} Original lemma 04XS: sources/stacks-stacks-more-morphisms/pinned-original-source.tex, source lines 2906--2906. Cited by the retained statement or proof but not part of the collected spans; supplied only as a navigation pointer. The argument does not depend on its content.\par\medskip
\ifdefined\hypertarget\hypertarget{sourcepointer-7}{}\fi
\noindent\textit{Source reference.} Original lemma 04XD: sources/stacks-stacks-more-morphisms/pinned-original-source.tex, source lines 2906--2906. Cited by the retained statement or proof but not part of the collected spans; supplied only as a navigation pointer. The argument does not depend on its content.\par\medskip
\begin{lemma}
\label{lemma-check-separated-dvr}
Let $f : \mathcal{X} \to \mathcal{Y}$ be a morphism of algebraic stacks. Assume
\begin{enumerate}
\item $\mathcal{Y}$ is locally Noetherian,
\item $f$ is locally of finite type and quasi-separated,
\item for every commutative diagram
$$
\xymatrix{
\Spec(K) \ar[r]_x \ar[d]_j & \mathcal{X} \ar[d]^f \\
\Spec(A) \ar[r]^y \ar@{-->}[ru] & \mathcal{Y}
}
$$
where $A$ is a discrete valuation ring and $K$ its fraction field
and any $2$-arrow $\gamma : y \circ j \to f \circ x$ the category
of dotted arrows (Morphisms of Stacks, Definition
\ref{stacks-morphisms-definition-fill-in-diagram})
is either empty or a setoid with exactly one isomorphism class.
\end{enumerate}
Then $f$ is separated.
\end{lemma}

\begin{proof}
To prove that $f$ is separated we have to show that
$\Delta : \mathcal{X} \to \mathcal{X} \times_\mathcal{Y} \mathcal{X}$
is proper. We already know that $\Delta$ is representable by
algebraic spaces, locally of finite type (Morphisms of Stacks,
Lemma \ref{stacks-morphisms-lemma-properties-diagonal}) and
quasi-compact and quasi-separated (by definition of $f$ being quasi-separated).
Choose a scheme $U$ and a surjective smooth morphism
$U \to \mathcal{X} \times_\mathcal{Y} \mathcal{X}$.
Set
$$
V = \mathcal{X} \times_{\Delta, \mathcal{X} \times_\mathcal{Y} \mathcal{X}} U
$$
It suffices to show that the morphism of algebraic spaces $V \to U$
is proper (Properties of Stacks, Lemma
\ref{stacks-properties-lemma-check-property-covering}).
Observe that $U$ is locally Noetherian
(use Morphisms of Stacks, Lemma
\ref{stacks-morphisms-lemma-locally-finite-type-locally-noetherian}
and the fact that $U \to \mathcal{Y}$ is locally of finite type)
and $V \to U$
is of finite type and quasi-separated (as the base change
of $\Delta$ and properties of $\Delta$ listed above). Applying
Cohomology of Spaces, Lemma \ref{spaces-cohomology-lemma-check-proper-dvr}
it suffices to show: Given a commutative diagram
$$
\xymatrix{
\Spec(K) \ar[r]_v \ar[d]_j &
V \ar[d]^g \ar[r] &
\mathcal{X} \ar[d]^\Delta \\
\Spec(A) \ar[r]^u \ar@{-->}[ru] \ar@{..>}[rru] &
U \ar[r] &
\mathcal{X} \times_\mathcal{Y} \mathcal{X}
}
$$
where $A$ is a discrete valuation ring and $K$ its fraction field,
there is a unique dashed arrow making the diagram commute.
By Morphisms of Stacks, Lemma
\ref{stacks-morphisms-lemma-cat-dotted-arrows-base-change}
the categories of dashed and dotted arrows are equivalent.
Assumption (3) implies there is a unique dotted arrow up to
isomorphism, see Morphisms of Stacks, Lemma
\ref{stacks-morphisms-lemma-helper-diagonal}. We conclude
there is a unique dashed arrow as desired.
\end{proof}

\end{document}
